\section{Week 7}

Last time we constructed cosets from an equivalence relation and looked at the set of all left cosets of subgroup \( N  \) of a group \( G  \), denoted by
\[  G / N = \{ gN : g \in G  \}.  \]
We showed that the operation on \( G / H  \) defined by
\[  (aH)(bH) = (ab) N \]
is well-defined when \( g n g^{-1} \in H  \) for all \( g \in G  \) and for all \( h \in N  \). This makes \( G  / N  \) a group. We will prove that this operation is associative, the rest is left to the reader. 

Let \( N \leq G  \) a group where "coset multiplication" is well-defined. Let \( aN, bN , cN \in G / N  \) \( (a,b,c \in G) \). Then 
\begin{align*}
    aN (bN cN) &= aN((bc)N) \\
               &= (a(bc)) N \\
               &= ((ab)c) N \\
               &= ((ab)N) cN \\
               &= (aN bN) cN.
\end{align*}

Similarly, we can show that \( 1N  \) is the identity of \( G / N   \) and show that 
\[  (aN)^{-1} = a^{-1}N \ \ \forall aN \in G / N.  \]

Indeed, for every \( g N \in G / N \) 
\[ g N =   (1_G g) N  =  (1_G N) \cdot (g N )    \]
and 
\[  g N = (g 1_G) N = (1_G N) \cdot (gN ). \]
Now, we have that since \( g \in G  \) and \( g^{-1} \in G  \), it follows that 
\[  {1}_{G} N = (g g^{-1})N = (g N) \cdot (g^{-1} N ) \]
and 
\[ {1}_{G} N = (g^{-1} g) N = (g^{-1} N ) \cdot (g N).  \]

This tells us that \( G /N  \) is a group where \( N  \) has the nice property of being "normal". This will be expanded in the following definition:

\begin{definition}[Normal Subgroups of \( G \)]
   A subgroup \( N  \) of \( G  \) is called \textbf{normal} in \( G  \) if every element of \( g  \) normalizes \( N  \). That is, \( N  \) is normal in \( G  \) if 
   \[  g N g^{-1} = N \ \ \forall g \in G.  \]
   If \( N  \) is a normal subgroup of \( G  \), then we write \( N \trianglelefteq G . \)
\end{definition}

\begin{theorem}[6]
    Let \( N  \) be a subgroup of \( G  \). The following are equivalent
    \begin{enumerate}
        \item[(1)] \( N \trianglelefteq G  \)
        \item[(2)] \( {N}_{G}(N) = G  \) 
        \item[(3)] \( g N = N g  \) for all \( g \in G  \)
        \item[(4)] The operation coset multiplication is well-defined.
        \item[(5)] \( g N g^{-1} \subset N  \) for all \( g \in G  \).
    \end{enumerate}
\end{theorem}

Why does (5) only need containment? It is because you can construct a bijective map between \( g N g^{-1}  \) and \( N  \).

\begin{eg}[Checking if \( \langle s \rangle  \) is normal in \( D_16 \)]
    We need to examine \( g s g^{-1} \) for an arbitrary \( g \in D_{16} \). Note that  
   \[  {D}_{16} = \{  1, r , \dots, r^{7}, s , sr, \dots, s r^{7} \}. \]
   Letting \( g = s^{i} r^{j} \) where \( i \in \{  0 , 1  \}  \) and \( j \in \{  0 , 1,2, \dots, 7 \}  \). Then we have 
   \begin{align*}
       gs g^{-1} &= (s^{i} r^{j}) s (s^{i} r^{j})^{-1} \tag{\( s^{i} = s^{-i} \)} \\
                 &= s^{i} r^{j} s r^{-j} s^{i} 
   \end{align*}
   Note that when \( i = 0  \), we have 
   \begin{align*}
       g s g^{-1} &= r^{j} s r^{-j} \\ 
                  &= r^{j} r^{j} s \\
                  &= r^{2j} s.
   \end{align*}
   When \( j = 1  \), this gives us \( g s g^{-1} = r^{2}s  \). But note that \( r^{2} s  \notin  \langle s \rangle \) because otherwise \( r^{2} \) is either the identity or \( s  \) which \( r^{2} \) is neither of these two which is a contradiction. Hence, \( g s g^{-1} \) cannot be in \( \langle s \rangle  \) and so it cannot be a normal subgroup of \( D_16 \).
\end{eg}

\begin{theorem}[Big!]
    A subgroup \( N \) of a group \( G  \) is normal if and only if it is the kernel of some homomorphism.
\end{theorem}
\begin{proof}
\( (\Longrightarrow) \) Suppose \( N \trianglelefteq G  \). Let's define the mapping \( \pi : G \to G /N  \) by \( \pi (g) = g N  \) for all \( g \in G  \). Notice that since \( N  \) is a normal subgroup of \( G  \), this map is well-defined. Let \( {g}_{1}, {g}_{2} \in G  \). Then 
\[  \pi({g}_{1} {g}_{2}) = ({g}_{1} {g}_{2}) N = ({g}_{1}N) ({g}_{2} N ) = \pi({g}_{1}) \pi({g}_{2}). \]
Hence, \( \pi  \) is a homomorphism. It remains to show that \( \text{ker}(\pi) = N  \). Indeed, we have that 
\begin{align*}
    \text{ker}(\pi) &= \{  g \in G : \pi(g) = 1N \}  \\
                    &= \{ g \in G : g N = 1 N  \}  \\
                    &= \{ g \in G : g \in 1N \}  \\
                    &= \{ g \in G : g \in N  \}  \\
                    &= N. 
\end{align*}
Hence, \( \text{ker}(\pi) = N  \).

\( (\Longleftarrow) \) This direction will be on homework.
\end{proof}

\begin{definition}[Natural Projection]
    Let \( N  \) be normal in \( G  \). The homomorphism \( \pi  \) from the previous theorem is called the \textbf{natural projection} (homomorphism) of \( G  \) onto \( G / N  \). If \( \overline{G} \leq G / N  \), the complete pre-image of \( \overline{H} \) is \( \pi^{-1}(\overline{H}) \).
\end{definition}

\begin{remark}
    If \( \overline{H} \subseteq  G / N  \), then \( N \leq \pi^{-1}(\overline{H}) \) since 
    \[  1 N \in \overline{H} \implies N = \text{ker}(\pi) = \pi^{-1}(1N) \subseteq  \pi^{-1}(\overline{H}). \]
\end{remark}

\begin{eg}[Complete Pre-image of \( Q_8 \)]
  \( \pi^{-1}(\overline{H}) = \{  1 , -1 , i , - i  \}  \) where   
  \begin{align*}
      \pi(i) &= i \langle - 1 \rangle   \\
      \pi(-1) &= - 1 \langle -1 \rangle = 1 \langle -1 \rangle
  \end{align*}
\end{eg}

\subsection{More on Normal Subgroups}

There are a lot of ways to see if a subgroup is normal.  

\begin{remark}[Important]
    Let \( G  \) be a group. Observe that \( \{  1  \}  \trianglelefteq G  \) and \( G \trianglelefteq G   \) and \( G / \{  1  \}  \cong G  \) and \( G / G \cong \{  1  \}  \).
\end{remark}

\begin{itemize}
    \item When \( G  \) is clearly an additive group, we denote left and right cosets \( g  + N  \) and \( N + g  \), respectively where \( N \leq G  \).
    \item When \( G  \) is abelian, every subgroup is normal; that is, \( N \leq G  \) where 
        \[  \forall n \in \N  \ \ \ gn g^{-1} = g g^{-1} n = n \in \N \ \ \forall g \in G.  \]
\end{itemize}

\subsection{More on Cosets/Lagrange's Theorem}

We move away from normal subgroups and just analyze subgroups.

\begin{theorem}[Lagrange's Theorem]
  If \( G  \) is a finite group, and \( H  \) is a subset of \( G  \), then   
  \[  | H  | \mid | G | \]
  and the number of left cosets of \( H  \) in \( G  \) is \( \frac{ | G |  }{  | H |  }  \).
\end{theorem}
The idea of the proof of this theorem is as follows: (see problem 18 and 19 from section 1.7 in the book)
\begin{proof}
Recall that left cosets form a partition of \( G  \); that is, 
\[  G = \bigcup_{  g \in G  }^{  }  g H.  \]
There is a bijection from \( H  \) to \( gH  \) and is defined by \( h \mapsto gh \) and so 
\[  | H |  = | g H  |. \]
At this point, we can conclude that \(  | G  |  = k | H |  \) where \( k  \) is the number of distinct left-cosets of \( H  \) in \( G  \). This gives us 
\[  k = \frac{ | G |   }{  | H |  }. \]
\end{proof}

\begin{definition}[Index of a Subgroup]
    If \( G  \) is a group (possibly infinite) and \( H  \) is a subgroup of \( G  \), then the number of distinct left cosets of \( H  \) in \( G  \) is called the \textbf{index of \( H  \) in \( G  \)} which is denoted by \( | G : H  |  \) (or \( [G : H]   \)). 
\end{definition}

\begin{corollary}[Corollary to Lagrange's Theorem]
   If \( G  \) is a finite group and \( x \in G  \), then \( | x  |  \mid | G  |  \).
\end{corollary}
\begin{proof}
    We proved that \( | x  |  = | \langle x \rangle  |   \) where \( \langle x \rangle \leq G  \). The rest follows from Lagrange's Theorem. 
\end{proof}

\begin{remark}
    For a finite group \( G  \) with \( H \leq G  \), \( | G : H  |  = \frac{ | G |  }{  | H  |  }  \).
\end{remark}

\begin{eg}
   Let \( G = \Z  \) and let \( H = 3 \Z  \), then \( |  \Z  : 3 \Z  |  \) is the number of left cosets of \( 3 \Z   \) in \( \Z  \). Recall that 
   \[  3 \Z =  \{  3x : x \in \Z  \}  \]
   where 
   \begin{align*}
       0 + 3 \Z &= 3 \Z  \\
       1 + 3 \Z &= \{  1 + 3x : x \in \Z  \}  \\
       2 + 3 \Z &= \{ 2 + 3x : x \in \Z  \}  \\
       3 + 3 \Z &= 0 + 3 \Z. 
   \end{align*}
   That is, 
   \[  | \Z : 3\Z  |  = 3 = | \Z / 3 \Z   |.  \]
\end{eg}

\begin{corollary}
   If \( G  \) is a group of prime order \( p  \), then \( G  \) is cyclic. 
\end{corollary}
\begin{proof}
Let \( x \in G  \) where \( x \neq 1  \). Then \( | x  |  \mid | G  |  \) since \( | G  |  = p  \) a prime. Hence, \( | x  | \in  \{  1, p \}  \). Since \( x \neq 1  \), we have \( | x  |  \neq 1  \) and so we would need to have \( | x  |  = p  \) and so \( \langle x \rangle = G  \).
\end{proof}

\begin{eg}
   A subgroup \( H  \) of a group \( G  \) with index \( | G : H  |  =  2  \), is normal. 
\end{eg}
\begin{proof}
Let \( g \in G - H  \), then \( gH \neq 1H  \). Since \( | G : H  |  = 2  \), there are two distinct left cosets of \( H  \) in \( G  \) and since one of them is \( 1  H  \), the other must be \( g H  \). Similarly, there are only two distinct right cosets of \( H  \) in \( G  \); namely, \( H1   \) and \( Hg  \). Since \( 1 H = H 1  \) and cosets form a partition of \( G  \), we have that \( gH = G - H =  H g  \). Hence, the left and right cosets of \( H  \) are equal and so \( H  \) is normal in \( G  \).
\end{proof} 

\begin{eg}[Normal Subgroups of is NOT transitive]
   Note that the relation "is a normal subgroup of" is not transitive. Let \( G = {D}_{8} \). Note that the \( | {D}_{8} |  = 8  \) and \( | \langle s \rangle |  = 2 \), \( | \langle s, r^{2} \rangle  |  = 4  \). Clearly, we have 
   \[  \langle s \rangle \subseteq  \langle s, r^{2} \rangle \subseteq  {D}_{8}. \]
   We have 
   \[  | {D}_{8} : \langle s, r^{2} \rangle  |  = 2  \]
   and \( | \langle s, r^{2} \rangle : \langle s \rangle  |  = 2  \). Thus, 
   \[  \langle s \rangle \trianglelefteq \langle s , r^{2} \rangle \trianglelefteq {D}_{8} \]
   but \( \langle s \rangle  \) is NOT normal in \( {D}_{8} \). Note that \( r s \neq sr  \) but \( rs = s r^{-1} \) where \( r^{2} s  \) is not \( 1  \) and not \( s  \).
\end{eg}

\begin{definition}[ ]
    Let \( H,K \leq G  \) where \(  G  \) is a group. Define
    \[  H K = \{  h k : h \in H, \ k \in K \}. \]
\end{definition}

\begin{theorem}[Big Theorem of this lecture]
   If \( H  \) and \( K  \) are finite subgroups of a group, then 
   \[  | H K |  = \frac{ | H  |  | K  |  }{  |  H \cap K  |  } \]
   where \( HK  \) need not be a group for the above to hold.
\end{theorem}
\begin{proof}
Notice that \( HK \) is the union of left cosets of \( K  \); that is,  
\[  HK = \bigcup_{ h \in H  }^{  }  h K.  \]
Since each coset of \( K  \) has \( | K  |  \) elements, we will count the number of distinct cosets in the above union. We know \( {h}_{1} k = {h}_{2} k  \) for \( {h}_{1}, {h}_{2} \in H  \) if and only if \( {h}_{2}^{-1} {h}_{1} \in  K  \). It follows that   
\begin{align*}
    {h}_{1} k = {h}_{2} k &\iff {h}_{2}^{-1} {h}_{1} \in K \cap H  \\
                          &\iff {h}_{1} (K \cap H) = {h}_{2} (K \cap H).
\end{align*}
Thus, the number of distinct cosets of the form \( h K  \) with \( h \in H  \) is the same as the number of distinct cosets of \( K \cap H  \) in \( H  \). Since we have that \( K \cap H  \) is a subgroup of \( H  \), this is \( \frac{ | H  |   }{  | K \cap H  |  } \). Therefore, \( HK  \) consists of \( \frac{ | H  |   }{  | H \cap K  |  }  \) distinct cosets of \( K  \) of which contains \( | K  |  \) elements, it follows that 
\[  | HK  |  = \frac{ | H  | | K  |   }{  | H \cap K  |  }.  \]
\end{proof}

\begin{eg}[An example where \( HK \) is not a subgroup]
    
\end{eg}




